\chapter{Zwartość i sklejanie trajektorii Morse'a: podstawy analityczne}
\label{ch:morse-analytic-v11}

\section{Cel i hipotezy twierdzeń}
W rozdziale~\ref{ch:morse-v10} skonstruowaliśmy różniczkę
Morse'a i użyliśmy kompaktyfikacji przestrzeni trajektorii.
Pełne konstrukcje geometryczne funkcji Morse'a,
lokalnych przepływów i uchwytów omówiliśmy też w
\href{https://github.com/dariusz-piekarz/mathematical_scipts/blob/biblioteka-matematyczna-iteracja-6/chapters/pl/differential_geometry_and_topology/18-teoria-morse-a-punkty-krytyczne-uchwyty-i-przeplyw.tex}{rozdziale
\emph{Teoria Morse'a: punkty krytyczne, uchwyty i przepływ}
skryptu \emph{Geometria różniczkowa}}
(kanoniczna etykieta: \texttt{ch:teoria-morsea}).
Nie należy jednak utożsamiać trzech odmiennych faktów:
(i) istnienia granic ciągów trajektorii, być może łamanych,
(ii) faktu, że każda trajektoria łamana jest granicą trajektorii
niełamanych, oraz (iii) zgodności znaków orientacji na
brzegu kompaktyfikacji. Pierwszy fakt wynika w dużej mierze
z ODE, zwartości i hipotezy Morse'a--Smale'a; drugi wymaga
argumentu sklejania, a trzeci reguły dla orientacji
przestrzeni przecięć. Rozdzielimy je i podamy dokładnie
hipotezy każdego kroku.

W całym rozdziale $M$ oznacza \emph{zamkniętą} (zwartą, bez brzegu)
gładką rozmaitość, $f:M\to\R$ funkcję Morse'a, $g$ metrykę
riemannowską, a $X=-\operatorname{grad}_g f$ jej ujemne pole
gradientowe. Przez $\Phi_t$ oznaczamy zupełny przepływ $X$.
Gdy potrzebujemy wymiarów przestrzeni moduli, dodatkowo zakładamy
warunek Morse'a--Smale'a. Nie wymagamy orientowalności $M$;
orientacje wybieramy oddzielnie na rozmaitościach niestabilnych.

\section{Równanie energii i granice trajektorii}
\begin{lemat}[Równanie energii i punkty graniczne]
\label{lem:morse-energy-analytic-v11}
Niech $M$ będzie zwartą gładką rozmaitością bez brzegu,
$f:M\to\R$ funkcją Morse'a, $g$ metryką riemannowską,
a $\gamma:\R\to M$ krzywą całkową pola
$X=-\operatorname{grad}_g f$.
Wówczas dla dowolnych $a,b\in\R$ z $a<b$ zachodzi
\[
 f(\gamma(a))-f(\gamma(b))=
 \int_{t\in[a,b]}\|\operatorname{grad}_g f(\gamma(t))\|_g^2\,\dd t.
\]
Ponadto zbiory graniczne $\omega(\gamma)$ i $\alpha(\gamma)$ są
niepustymi, spójnymi podzbiorami $\operatorname{Crit}(f)$;
w konsekwencji są pojedynczymi punktami krytycznymi.
\end{lemat}
\begin{proof}
Pierwsza tożsamość wynika z
$\tfrac{\dd}{\dd t}f(\gamma(t))=-\|\nabla f(\gamma(t))\|^2$.
Monotoniczność i zwartość $M$ dają skończone granice
$\ell_+=\lim_{t\to+\infty}f(\gamma(t))$ oraz
$\ell_-=\lim_{t\to-\infty}f(\gamma(t))$.
Zbiór $\omega(\gamma)$ jest przecięciem malejącej rodziny
niepustych, zwartych i spójnych zbiorów
$\overline{\gamma([T,\infty))}$, więc jest niepusty i spójny.
Jest niezmienniczy względem przepływu. Jeżeli punkt $y$ tego
zbioru byłby niekrytyczny, to
$\tfrac{\dd}{\dd t}f(\Phi_t(y))|_{t=0}<0$.
Tymczasem na całym zbiorze $\omega(\gamma)$ funkcja $f$
ma stałą wartość $\ell_+$, a niezmienniczość daje sprzeczność.
Zatem $\omega(\gamma)\subseteq\operatorname{Crit}(f)$.
Zbiór krytyczny jest skończony, a spójny podzbiór zbioru
skończonego jest jednopunktowy. Rozumowanie dla $\alpha(\gamma)$
jest analogiczne, po zmianie kierunku czasu. Ponieważ zbiory
graniczne są jednopunktowe, mamy faktyczną zbieżność krzywej
przy $t\to\pm\infty$.
\end{proof}

\begin{lemat}[Jednostajne ograniczenie czasu z dala od zer]
\label{lem:morse-transit-v11}
Niech $M$ będzie zwartą gładką rozmaitością bez brzegu,
$f:M\to\R$ funkcją Morse'a, a $g$ metryką riemannowską.
Niech $U\subseteq M$ będzie otwartym otoczeniem zbioru
$\operatorname{Crit}(f)$ oraz $K=M\setminus U$.
Jeżeli $K\ne\varnothing$, to istnieje $m>0$ takie, że
$\|\nabla f(x)\|_g\ge m$ dla $x\in K$.
Łączny czas, jaki dowolna niestała trajektoria gradientowa
spędza w $K$, nie przekracza
$\bigl(\max_M f-\min_M f\bigr)/m^2$.
\end{lemat}
\begin{proof}
Zwartość $K$ i brak punktów krytycznych dają
$m=\min_K\|\nabla f\|_g>0$.
Dla dowolnego mierzalnego podzbioru $A\subset\R$, dla którego
$\gamma(A)\subseteq K$, równanie energii daje
$m^2\lambda_1(A)\le\int_{t\in A}\|\nabla f(\gamma(t))\|_g^2\,\dd t
\le\max_Mf-\min_Mf$.
Tutaj $\lambda_1$ oznacza jednowymiarową miarę Lebesgue'a.
W szczególności nierówność zachodzi dla całego zbioru czasów,
w których krzywa przebywa w $K$.
\end{proof}

\section{Granice ciągów rozwiązań i trajektorie łamane}
\begin{definicja}[Trajektoria łamana]
\label{def:broken-morse-v11}
Niech $M$ będzie zwartą gładką rozmaitością bez brzegu,
$f:M\to\R$ funkcją Morse'a, $g$ metryką riemannowską,
a $X=-\operatorname{grad}_g f$.
Dla $p,q\in\operatorname{Crit}(f)$ \emph{trajektorią łamaną}
od $p$ do $q$ nazywamy niepusty skończony ciąg
$([\gamma_1],\ldots,[\gamma_\ell])$, $\ell\ge1$,
gdzie każda $\gamma_j:\R\to M$ jest niestałą krzywą całkową $X$,
$\gamma_1(-\infty)=p$, $\gamma_\ell(+\infty)=q$ oraz
$\gamma_j(+\infty)=\gamma_{j+1}(-\infty)$.
Nawias oznacza klasę z dokładnością do translacji parametru czasu.
\end{definicja}

\begin{twierdzenie}[Zwartość z dokładnością do łamania]
\label{thm:morse-broken-compact-v11}
Niech $M$ będzie zwartą gładką rozmaitością bez brzegu,
$f:M\to\R$ funkcją Morse'a, a $g$ metryką riemannowską.
Niech $p,q\in\operatorname{Crit}(f)$, $p\ne q$, i niech
$([\gamma_j])_{j\in\N}$ będzie ciągiem w przestrzeni
niesparametryzowanych trajektorii
$\mathcal M(p,q)=(W^u(p)\cap W^s(q))/\R$
ujemnego przepływu gradientowego.
Istnieje podciąg oraz trajektoria łamana
$([\eta_1],\ldots,[\eta_\ell])$ od $p$ do $q$ takie, że
każdą $\eta_i$ można otrzymać jako granicę $C^\infty$
na zwartych przedziałach odpowiednio \emph{przesuniętych w czasie}
krzywych $\gamma_j$. Ponadto wartości $f$ na kolejnych
punktach krytycznych są ściśle malejące, więc
$\ell\le\#\operatorname{Crit}(f)-1$.
Jeśli $(f,g)$ jest parą Morse'a--Smale'a, to
$\ell\le\lambda(p)-\lambda(q)$.
\end{twierdzenie}
\begin{proof}
Zbiór punktów krytycznych jest skończony. Wybierzmy
skończony zbiór regularnych wartości $c_1>\cdots>c_N$
tak, aby między każdymi dwiema kolejnymi różnymi wartościami
krytycznymi leżało jedno $c_i$ (niektóre $c_i$ mogą leżeć
poza przedziałem $(f(q),f(p))$ i wtedy je pomijamy).
Aby kontrolować \emph{także końce} trajektorii, wybieramy
odpowiednio blisko wartości $f(p)$ dodatkowy regularny
poziom $c_+$, położony poniżej $f(p)$, a blisko $f(q)$
poziom $c_-$, położony powyżej $f(q)$.
Można je wybrać tak, by każdy punkt przecięcia trajektorii
wychodzącej z $p$ z $f^{-1}(c_+)$ należał do zwartej
sfery przekrojowej lokalnej rozmaitości niestabilnej
$W^u_{\mathrm{loc}}(p)$, a analogicznie dla $c_-$
i $W^s_{\mathrm{loc}}(q)$. Rzeczywiście w lokalnych
współrzędnych stabilnych i niestabilnych funkcja $f$
ma na tych rozmaitościach odpowiednio niezdegenerowane
maksimum i minimum; poziomy blisko wartości krytycznej
są zwartymi sferami, a na brzegu wybranego małego
otoczenia spadek (względnie wzrost w czasie wstecznym)
ma dodatnie ograniczenie od dołu. Dlatego trajektoria
musi przeciąć taki poziom przed opuszczeniem otoczenia.
Ten wybór wyklucza utratę końcowego punktu $p$ lub $q$
na rzecz innego punktu krytycznego o tej samej wartości.

Na każdej niestałej trajektorii od $p$ do $q$ funkcja
$f\circ\gamma_j$ jest ściśle malejąca i przyjmuje każdą
wartość $c_i\in(f(q),f(p))$ dokładnie raz; do rodziny
przekrojów włączamy także $c_+$ i $c_-$, jeżeli nie były
już wybrane. Niech
$x_{j,i}$ będzie odpowiednim punktem przecięcia krzywej
z zwartą regularną hiperpłaszczyzną $f^{-1}(c_i)$.
Po wybraniu podciągu wszystkie $x_{j,i}$ zbiegają do pewnych
$x_i\in f^{-1}(c_i)$.

Każda krzywa $t\mapsto\Phi_t(x_{j,i})$ zbiega na zwartych
przedziałach do $\eta_i(t):=\Phi_t(x_i)$, nawet w topologii
$C^\infty$, ponieważ przepływ jest gładki, a $M$ zwarta.
Granice w obu kierunkach czasu są punktami krytycznymi
na mocy Lematu~\ref{lem:morse-energy-analytic-v11}.
Zachodzą dwa przypadki dla sąsiednich przekrojów $c_i,c_{i+1}$:
czas przejścia krzywej $\gamma_j$ między nimi ma podciąg
ograniczony albo dąży do $+\infty$.
W pierwszym przypadku $x_i$ oraz $x_{i+1}$ należą do tej
samej trajektorii granicznej, gdyż przejście jest realizowane
przepływem przez skończony czas.
W drugim przypadku graniczne odcinki są rozdzielone
punktem krytycznym: gdyby graniczna trajektoria nie
zbliżała się do krytycznego punktu w obszarze między
poziomami, Lemat~\ref{lem:morse-transit-v11}, użyty do
tego obszaru z usuniętymi małymi otoczeniami punktów
krytycznych, ograniczałby czasy przejścia.
Ponieważ przepływ nie może zwiększać $f$, a odstępy czasu
między dwoma sąsiednimi przekrojami są jedynym miejscem
możliwej utraty zwartości, przy sklejaniu
sąsiednich granicznych odcinków ich końce muszą być tym
samym punktem krytycznym; dwie różne krytyczne wartości
pośrednie zostałyby rozdzielone przez dodatkowy poziom
$c_i$. W rezultacie otrzymujemy skończony ciąg trajektorii
łączących $p$ z $q$.

Każdy nietrywialny odcinek powoduje ścisły spadek $f$,
co daje ograniczenie liczby odcinków przez liczbę
punktów krytycznych minus jeden.
Gdy para jest Morse'a--Smale'a, z wymiaru
$\dim\mathcal M(a,b)=\lambda(a)-\lambda(b)-1\ge0$
wynika $\lambda(a)>\lambda(b)$ dla każdej niepustej
przestrzeni trajektorii $\mathcal M(a,b)$.
Sumując różnice indeksów otrzymujemy
$\ell\le\lambda(p)-\lambda(q)$.
\end{proof}

\begin{wniosek}[Skończoność trajektorii o różnicy indeksów jeden]
\label{cor:morse-rigid-finite-v11}
Niech $M$ będzie zwartą gładką rozmaitością bez brzegu,
$(f,g)$ parą Morse'a--Smale'a na $M$ i niech
$p,q\in\operatorname{Crit}(f)$.
Jeśli $\lambda(p)-\lambda(q)=1$, to
$\mathcal M(p,q)$ jest skończonym zbiorem.
\end{wniosek}
\begin{proof}
Z twierdzenia wynika, że każdy ciąg w $\mathcal M(p,q)$
ma podciąg zbieżny z dokładnością do łamania.
Dwa lub więcej niezerowych odcinków wymagałoby spadku indeksu
co najmniej dwa, więc granica pozostaje niełamana.
Wybór przekroju z regularnym poziomem $f=c$ identyfikuje
zbieżność modulo przesunięcia czasu ze zwykłą zbieżnością
punktów na przekroju. Zatem $\mathcal M(p,q)$ jest zwarta.
Z warunku Morse'a--Smale'a jest zerowymiarową rozmaitością,
czyli przestrzenią dyskretną; zwarta przestrzeń dyskretna
jest skończona.
\end{proof}

\begin{uwaga}[Co uzyskaliśmy bez twierdzenia o sklejaniu?]
Zwartość z dokładnością do łamania i skończoność
$\mathcal M(p,q)$ przy różnicy indeksów jeden są dowiedzione
bez używania konstrukcji gluing (przy przyjęciu
standardowego twierdzenia o rozmaitościach stabilnych
i gładkiej teorii ODE). Natomiast z tych dwóch
faktów \emph{nie wynika}, że każda trajektoria łamana
jest końcem jednowymiarowej przestrzeni trajektorii.
Dopiero twierdzenie o sklejaniu nadaje temu końcowi
strukturę przedziału i pozwala liczyć znaki brzegowe.
\end{uwaga}

\section{Analiza w otoczeniu punktu hiperbolicznego}
\begin{lemat}[Wykładniczy zanik na rozmaitości stabilnej]
\label{lem:morse-exponential-v11}
Niech $M$ będzie zwartą gładką rozmaitością bez brzegu,
$f:M\to\R$ funkcją Morse'a, $g$ metryką riemannowską,
a $q\in\operatorname{Crit}(f)$ jej punktem krytycznym.
Istnieją otoczenie $V$ punktu $q$ oraz liczby $C,\mu>0$
takie, że każda trajektoria z lokalnej rozmaitości stabilnej,
która od czasu $0$ pozostaje w $V$, spełnia
$\operatorname{dist}_g(\gamma(t),q)\le
C e^{-\mu t}\operatorname{dist}_g(\gamma(0),q)$ dla $t\ge0$.
Analogiczna nierówność dla $t\le0$ zachodzi na lokalnej
rozmaitości niestabilnej.
\end{lemat}
\begin{proof}
W lokalnych współrzędnych przystosowanych do rozmaitości
stabilnej równanie na tej rozmaitości ma postać
$\dot v=A_sv+R(v)$, gdzie wszystkie wartości własne
$A_s$ mają ujemne części rzeczywiste oraz $R(0)=0$,
$\D R(0)=0$. Istnieje dodatnio określona macierz $P$
spełniająca równanie Lapunowa
$A_s^\top P+PA_s=-I$; można ją otrzymać przez
$P=\int_{t\in[0,\infty)}e^{tA_s^\top}e^{tA_s}\,\dd t$.
Dla $Q(v)=v^\top Pv$ mamy
$\dot Q=-\|v\|^2+2v^\top PR(v)$.
Po zmniejszeniu otoczenia składnik ostatni ma moduł
co najwyżej $\frac12\|v\|^2$, zatem
$\dot Q\le-\mu_0 Q$ dla pewnego $\mu_0>0$.
Równoważność norm i nierówność Gronwalla kończą dowód.
Lokalne równanie na rozmaitości niestabilnej rozpatrujemy
po odwróceniu kierunku czasu. W dowodzie użyliśmy istnienia
gładkich lokalnych rozmaitości stabilnych z twierdzenia
o rozmaitościach stabilnych.
\end{proof}

\begin{lemat}[Jawny prawy odwrotny modelu jednowymiarowego]
\label{lem:morse-scalar-green}
Niech $\lambda\in\R\setminus\{0\}$,
a $h:\R\to\R$ będzie ciągłą funkcją ograniczoną.
Równanie
\[
 u'(t)+\lambda u(t)=h(t),\qquad t\in\R,
\]
ma dokładnie jedno rozwiązanie $u\in C^1(\R)$
ograniczone na całej prostej. Jest ono dane przez
\[
 (Q_\lambda h)(t)=
 \begin{cases}
 \displaystyle\int_{s\in(-\infty,t]}
           e^{-\lambda(t-s)}h(s)\,\dd s,&\lambda>0,\\[6pt]
 \displaystyle-\int_{s\in[t,\infty)}
           e^{-\lambda(t-s)}h(s)\,\dd s,&\lambda<0.
 \end{cases}
\]
Ponadto $\|Q_\lambda h\|_\infty
\le|\lambda|^{-1}\|h\|_\infty$, a oszacowanie
nie zależy od przesunięcia argumentu funkcji $h$.
\end{lemat}
\begin{proof}
Dla $\lambda>0$ moduł całkowanej funkcji jest
ograniczony przez $e^{-\lambda(t-s)}\|h\|_\infty$,
którego całka po $s\le t$ wynosi $1/\lambda$.
Daje to istnienie całki i oszacowanie normy.
Różniczkując po $t$ na mocy wzoru Leibniza,
otrzymujemy $u'=h-\lambda u$.
Gdy $\lambda<0$, całkujemy od $t$ ku $+\infty$:
wówczas $t-s\le0$ i
$e^{-\lambda(t-s)}=e^{-|\lambda|(s-t)}$,
a całka z jądra wynosi $1/|\lambda|$.
Pochodna drugiej formuły również jest
$u'=h-\lambda u$.

Różnica dwóch ograniczonych rozwiązań
spełnia równanie jednorodne, więc ma postać
$C e^{-\lambda t}$.
Dla $\lambda>0$ jest nieograniczona
przy $t\to-\infty$, a dla $\lambda<0$
przy $t\to+\infty$, chyba że $C=0$.
Stąd jednoznaczność. Wreszcie zmiana
$h(t)$ na $h(t-\tau)$ tylko przesuwa
przedstawione całki; stała oszacowania
pozostaje ta sama.
\end{proof}

\begin{uwaga}[Dlaczego model skalarny nie domyka ogólnego sklejania]
Lemat~\ref{lem:morse-scalar-green} dostarcza
rzeczywistego operatora Greena i jednostajnego oszacowania
dla stałego operatora hiperbolicznego.
W twierdzeniu o sklejaniu operator
$\mathcal D_T$ ma jednak \emph{zmienne współczynniki}
wzdłuż obu trajektorii, a ponadto
kierunek translacji czasu i kierunek długości
szyjki wymagają normalizacji.
Nie wolno więc zastąpić lematu o jednostajnym
prawym odwrotnym dla $\mathcal D_T$
samym powyższym rachunkiem skalarnym.
\end{uwaga}

\begin{uwaga}[Dlaczego to nie wystarcza do sklejania?]
Zanik rozwiązań ustalonych na $W^s(q)$ i $W^u(q)$ nie daje
jeszcze rozwiązania, które przechodzi \emph{blisko} $q$,
ale nie należy do żadnej z tych rozmaitości. Aby je
skonstruować, trzeba kontrolować równanie dla małych
odchyleń w obu kierunkach czasu jednocześnie.
Kluczowe są oszacowania dla operatora
$\mathcal D_\gamma=\tfrac{\dd}{\dd t}-\D X(\gamma(t))$
w odpowiednich ważonych przestrzeniach funkcyjnych.
\end{uwaga}

\section{Sklejanie i orientacje jednowymiarowych przestrzeni moduli}
\begin{twierdzenie}[Model gluing w różnicy indeksów dwa]
\label{thm:morse-gluing-v11}
Niech $M$ będzie zwartą gładką rozmaitością bez brzegu,
$(f,g)$ parą Morse'a--Smale'a na $M$, a punkty
$p,q,r\in\operatorname{Crit}(f)$ niech spełniają
$\lambda(p)-\lambda(q)=\lambda(q)-\lambda(r)=1$.
Dla każdej pary trajektorii
$([\gamma_-],[\gamma_+])\in
\mathcal M(p,q)\times\mathcal M(q,r)$
istnieją $T_0>0$ oraz gładkie odwzorowanie
\[
\operatorname{Glue}_{\gamma_-,\gamma_+}:
(T_0,\infty)\longrightarrow\mathcal M(p,r),
\]
którego wartości zbiegają do danej trajektorii łamanej
w sensie Twierdzenia~\ref{thm:morse-broken-compact-v11}.
Dla dostatecznie dużego $T$ odwzorowanie jest lokalnie różnowartościowe,
a każda dostatecznie bliska trajektoria niełamana należy do
jej obrazu. Parametr $e^{-T}$ opisuje topologicznie zbliżanie się do
brzegu; dla struktury gładkiego kołnierza należy dobrać
parametr sklejania zgodny z uzyskanymi oszacowaniami.
Nie twierdzimy bez dodatkowego argumentu, że sama funkcja
$e^{-T}$ jest zawsze gładką współrzędną na brzegu.
\end{twierdzenie}
\begin{proof}[Dowód warunkowy: wymagany lemat analityczny]
Przesuwamy parametryzacje $\gamma_-$ i $\gamma_+$ tak,
aby na wybranych przekrojach poprzecznych do przepływu
przy $q$ miały ustalone czasy przejścia. Wstawiamy między
ich końcowe fragmenty odcinek długości w przybliżeniu $2T$.
Wykładniczy zanik z Lematu~\ref{lem:morse-exponential-v11}
pozwala sklejać krzywe za pomocą gładkich funkcji odcięcia
w przybliżone rozwiązanie $u_T$, którego defekt
$\dot u_T-X(u_T)$ jest rzędu $O(e^{-\mu T})$.

Aby uzyskać rzeczywiste rozwiązanie, rozważamy w lokalnej
trywializacji operator
$\mathcal D_T=\tfrac{\dd}{\dd t}-\D X(u_T(t))$.
Warunek Morse'a--Smale'a daje surjektywność odpowiednich
operatorów linearyzacji wzdłuż $\gamma_-$ i $\gamma_+$.
Niezbędny jest następujący \emph{lemat analityczny o prawym
odwrotnym}: po wybraniu wag $0<\delta<\mu$ mniejszych od luk
spektralnych hesjanu w punktach końcowych i po
usunięciu kierunku translacji czasu (w odpowiednim
przekroju afinicznym) rozpatrujemy operator między
ważonymi przestrzeniami Sobolewa $W^{1,2}_\delta$
i $L^2_\delta$. Dopiero po wydzieleniu kierunku
parametru sklejania wymagamy istnienia prawego odwrotnego
$Q_T$ i stałej $C$ niezależnej od dużego $T$, takich że
$\mathcal D_TQ_T=\id$, $\|Q_T\|\le C$, w odpowiednich
normach ważonych. Buduje się ją przez sklejanie prawych
odwrotności dla $\gamma_-$ i $\gamma_+$ oraz poprawienie
błędu operatorem Neumanna. Ten lemat wymaga kontroli
jednowymiarowego kierunku związanego z długością szyjki;
nie wynika automatycznie z oddzielnej surjektywności
operatorów dla $\gamma_-$ i $\gamma_+$.

Pisząc poprawkę jako $u_T+\xi$, sprowadzamy równanie do
$\xi=-Q_T\bigl(\dot u_T-X(u_T)+N_T(\xi)\bigr)$,
gdzie $\|N_T(\xi)-N_T(\eta)\|\le
C(\|\xi\|+\|\eta\|)\|\xi-\eta\|$ dla małych norm.
Jednostajny szacunek na $Q_T$ i zanik defektu zapewniają,
że prawa strona jest kontrakcją w kuli promienia
$O(e^{-\mu T})$. Stąd otrzymujemy rodzinę prawdziwych
trajektorii i jej gładką zależność od $T$.
Argument \emph{jednoznaczności} w kołnierzu opiera się na
tej samej konstrukcji odwrotności oraz na analizie
przekrojów: każda bliska trajektoria ma jednoznaczną
normalizację czasu i długość szyjki.

Zatem główne równanie korekcyjne i oszacowania są opisane
jawnie, ale sam lemat o jednostajnym prawym odwrotnym
przyjmujemy tutaj jako standardowe twierdzenie analityczne.
Nie należy uważać powyższego schematu za jego pełny
dowód niezależny.
\end{proof}

\begin{przyklad}[Sklejanie trajektorii na torusie: model obliczony dokładnie]
\label{prz:morse-torus-gluing-explicit}
Rozważmy torus $T^2=(\R/2\pi\Z)^2$ z płaską
metryką produktową i funkcją
\[
 f(\theta,\varphi)=\cos\theta+\cos\varphi.
\]
Jej ujemne pole gradientowe ma równanie
\[
 \dot\theta=\sin\theta,\qquad
 \dot\varphi=\sin\varphi.
\]
Punkt $p=(0,0)$ jest maksimum o indeksie $2$,
$r=(\pi,\pi)$ minimum o indeksie $0$,
a $q_\theta=(\pi,0)$ i $q_\varphi=(0,\pi)$
są siodłami o indeksie $1$.
To para Morse'a--Smale'a:
rozmaitości stabilne i niestabilne są produktami
odpowiednich łuków okręgów, a ich
przecięcia są transwersalne.
Obydwa siodła mają tę samą wartość krytyczną $0$,
co nie przeszkadza rozpatrywaniu trajektorii łamanych.

W jednym wymiarze równanie $\dot\theta=\sin\theta$
możemy rozwiązać przez podstawienie
$z=\tan(\theta/2)$.
Ponieważ $\dot z=z$, wszystkie niestałe
trajektorie od $0$ do $\pi$ modulo $2\pi$
mają postać
\[
 \theta_{\varepsilon,a}(t)=
 2\varepsilon\arctan(e^{t-a}),\qquad
 \varepsilon\in\{-1,1\},\quad a\in\R.
\]
Dwa znaki odpowiadają dwóm łukom okręgu.
Każda trajektoria od $p$ do $r$ ma więc
jednoznaczny opis
\[
 \gamma_{\varepsilon,\delta,a,b}(t)=
 \bigl(2\varepsilon\arctan(e^{t-a}),
       2\delta\arctan(e^{t-b})\bigr),
 \qquad \varepsilon,\delta\in\{-1,1\}.
\]
W obu współrzędnych musi zachodzić przejście
od $0$ do $\pi$; rozwiązanie ze stałą
współrzędną nie może łączyć $p$ z $r$.
Przesunięcie czasu o $\tau$ zmienia
$(a,b)$ na $(a-\tau,b-\tau)$.
Dlatego niezmiennikiem klasy trajektorii
niesparametryzowanej jest
$s=b-a\in\R$, a dla każdej pary
znaków $(\varepsilon,\delta)$ dostajemy
\emph{jedną pełną składową}
$\mathcal M(p,r)$ dyfeomorficzną z $\R$.
Są dokładnie cztery takie składowe.

Ustalmy reprezentanta z $a=-s/2$, $b=s/2$.
Dla $s\to+\infty$ współrzędna $\theta$
zmienia się dużo wcześniej niż $\varphi$;
po niezależnym przesunięciu czasu na
obu fragmentach granicą jest para trajektorii
$p\to q_\theta\to r$.
Dla $s\to-\infty$ otrzymujemy
$p\to q_\varphi\to r$.
W każdym przypadku zachowujemy znaki
$\varepsilon,\delta$, więc są dokładnie
cztery pary trajektorii łamanych
przez każde z siodeł, po jednej dla
każdej składowej odpowiedniej przestrzeni moduli.

W tym \emph{konkretnym modelu}
kompaktyfikacja każdej kopii $\R_s$
jest domkniętym przedziałem $[-\infty,+\infty]$.
Współrzędną kołnierzową przy prawym końcu
można wybrać jako $\rho_+=e^{-s}$,
a przy lewym jako $\rho_-=e^{s}$.
Nie wystarczy tylko nazwać tego parametru:
sprawdźmy gładkość na samych formułach.
Dla $s\to+\infty$ podstawmy $\rho=e^{-s}$.
W przekroju czasowym przy pierwszym przejściu,
$t=-s/2+u$, mamy
\[
 \theta=2\varepsilon\arctan(e^u),\qquad
 \varphi=2\delta\arctan(\rho e^u).
\]
W przekroju przy drugim przejściu,
$t=s/2+u$, mamy z kolei
\[
 \theta=\varepsilon\pi-2\varepsilon
           \arctan(\rho e^{-u}),\qquad
 \varphi=2\delta\arctan(e^u).
\]
Dla $u$ w dowolnym przedziale zwartym
obie pary formuł zależą gładko od
$\rho\in[0,\varepsilon_0)$, również przy $\rho=0$;
są to dokładnie dwie składowe granicznej trajektorii
łamanej przez $q_\theta$.
Przy $s\to-\infty$ zamieniamy role
$\theta$ i $\varphi$ i otrzymujemy tę samą
własność dla $\rho_-=e^s$.
Pokazuje to jawnie istnienie gładkich
jednostronnych parametrów sklejania
dla \emph{tej jednej funkcji i metryki}.
Nie trzeba tutaj postulować
abstrakcyjnego operatora $Q_T$,
aby skonstruować tę kompaktyfikację.

Każdy z czterech zorientowanych przedziałów
ma przeciwne znaki dwóch punktów brzegowych.
Po wybraniu orientacji niestabilnych
rozmaitości również odpowiednie dwie
trajektorie od siodła do minimum
znoszą się znakami, podobnie jak
dwie trajektorie od maksimum do siodła.
Stąd dla tego modelu
$\partial_2=0$ i $\partial_1=0$,
co zgadza się z
$H_*(T^2;\Z)\cong(\Z,\Z^2,\Z)$
w stopniach $0,1,2$.
Przykład jest \emph{pełnym, szczególnym dowodem sklejania},
ale nie zastępuje ogólnego lematu
o prawym odwrotnym dla dowolnej pary Morse'a--Smale'a.
\end{przyklad}

\begin{figure}[!htbp]
\centering
\begin{tikzpicture}[>=Stealth,font=\small]
 \draw[->,thick] (-4,0)--(4.15,0)
       node[right] {$s=b-a$};
 \fill[manifoldblue] (-3.55,0) circle (2.5pt);
 \fill[manifoldblue] (3.55,0) circle (2.5pt);
 \node[above,align=center] at (-3.55,.18)
       {$p\to q_\varphi\to r$};
 \node[above,align=center] at (3.55,.18)
       {$p\to q_\theta\to r$};
 \node[below=7pt] at (0,-.15)
       {jedna z czterech składowych $\mathcal M(p,r)\cong\R$};
 \node[below=26pt,manifoldred] at (-3.4,-.15)
       {$\rho_-=e^{s}\to0$};
 \node[below=26pt,manifoldred] at (3.4,-.15)
       {$\rho_+=e^{-s}\to0$};
 \node[above=33pt] at (0,0)
       {$\overline{\mathcal M}(p,r)\cong[-\infty,+\infty]$};
\end{tikzpicture}
\caption{Jawna kompaktyfikacja składowej przestrzeni
trajektorii na torusie. Parametr $s$ mierzy
odstęp czasów przejścia obu współrzędnych;
dwa krańce są różnymi trajektoriami łamanymi.}
\label{fig:morse-torus-gluing-explicit}
\end{figure}

\begin{center}
\begin{tikzpicture}[>=Stealth,font=\small,x=1cm,y=1cm]
\coordinate (P) at (-4,0.5); \coordinate (Q) at (0,0); \coordinate (R) at (4,-0.5);
\fill[manifoldblue] (P) circle (2pt);
\fill[manifoldblue] (Q) circle (2pt);
\fill[manifoldblue] (R) circle (2pt);
\draw[manifoldblue,very thick,-{Stealth[length=2mm]}] (P) .. controls (-2.7,.75) and (-1.6,.15) .. (-.25,.05);
\draw[manifoldblue,very thick,-{Stealth[length=2mm]}] (.25,-.06) .. controls (1.7,-.45) and (2.8,-.85) .. (R);
\node[above] at (P) {$p$};\node[above] at (Q) {$q$};\node[above] at (R) {$r$};
\draw[densely dotted,manifoldgold,thick] (-.9,-.65) rectangle (.9,.65);
\node[below] at (0,-.78) {otoczenie $q$; czas szyjki $2T$};
\node[above] at (-2.25,.55) {$\gamma_-$};
\node[above] at (2.3,-.25) {$\gamma_+$};
\end{tikzpicture}
\par\smallskip\small Trajektoria łamana $p\to q\to r$.
Sklejanie zastępuje dwa odcinki jedną trajektorią
przechodzącą blisko $q$ przez długi czas.
\end{center}

\begin{stwierdzenie}[Orientacje i różniczka kwadratowa]
\label{prop:morse-d2-analytic-v11}
Przy założeniach Twierdzenia~\ref{thm:morse-gluing-v11}
i zgodnych orientacjach rozmaitości niestabilnych,
zorientowane końce jednowymiarowej kompaktyfikacji
$\overline{\mathcal M}(p,r)$ są oznaczone iloczynami
$\varepsilon(\gamma_-)\varepsilon(\gamma_+)$
z jednym, wspólnym znakiem konwencji brzegowej.
W konsekwencji dla kompleksu Morse'a $\partial^2=0$,
o ile zastosujemy twierdzenie o zgodności orientacji
przy sklejaniu (niezależnie przywołane poniżej).
\end{stwierdzenie}
\begin{proof}
W pobliżu końca gluing orientacja jednowymiarowego
moduli jest orientacją parametru długości szyjki,
a znak pozostałych czynników pochodzi z orientacji
transwersalnych przecięć $W^u(p)$ z $W^s(q)$ i
$W^u(q)$ z $W^s(r)$.
Zgodność znaku tej operacji z regułą brzegu należy do
standardowej \emph{orientacyjnej części twierdzenia
o sklejaniu}. Po jej zastosowaniu każdy punkt
\(\mathcal M(p,q)\times\mathcal M(q,r)\)
wnosi do współczynnika $r$ w $\partial^2p$ swój
zorientowany znak brzegowy (pomnożony przez stały
znak globalny). Zorientowana zwarta
jednowymiarowa rozmaitość ma sumę znaków punktów
brzegowych równą zeru. Stąd każdy współczynnik
$\partial^2p$ znika. Pełny dowód reguły orientacji
wymaga porównania orientacji wyznaczników operatorów
linearyzacji i pozostaje odrębnym krokiem od
samej zwartości ODE.
\end{proof}

\section{Dokładny bilans: co jest, a czego jeszcze nie dowiedziono?}
Dzięki Lematowi~\ref{lem:morse-energy-analytic-v11},
Twierdzeniu~\ref{thm:morse-broken-compact-v11} i
Wnioskowi~\ref{cor:morse-rigid-finite-v11} mamy już
pełne uzasadnienie skończoności liczby sztywnych trajektorii,
przy użyciu jedynie standardowej teorii gładkich przepływów
i rozmaitości stabilnych. Dla jednowymiarowych moduli
pozostały dwa jasno nazwane narzędzia analityczne:
jednostajne prawe odwrotności dla operatorów gluing i
zgodność orientacji na brzegu. Wnioski homologiczne są
więc dowiedzione \emph{względem tych twierdzeń pomocniczych},
a nie przez sam rachunek energii. Jest to precyzyjniejszy
status niż stwierdzenie, że ``kompaktyfikacja wynika
ze zwartości rozwiązań ODE''.



